\scanpage{371}
\mspage{188}

Let us recall that, via the choice of $U_{g,\nu}$ (differentiable), we had been able to construct by transcendental means a $\widetilde M_{g,\nu\,\mathbb C}$\inserted{ and a $\unsure{\widetilde{\mathcal U}_{g,\nu\,\mathbb C}|\mathcal U_{g,\nu\,\mathbb C}}=\unsure{\widetilde M_{g,\nu+1\,\mathbb C}|M_{g,\nu+1\,\mathbb C}}$}, hence a $\widetilde M_{g,\nu\,\mathbb Q}$, \deleted{hence \unsure{a|a}}\inserted{and a} $\widetilde{\mathcal U}_{g,\nu\,\mathbb Q}$, hence a $\pi_1(\mathcal U_{g,\nu,\mathbb Q})$, with a filtration in three \unsure{steps}, whose factors are respectively canonically isomorphic to
\[
\widehat{\pi}_{g,\nu},\quad
\unsure{\widehat{\mathcal T(\pi_{g,\nu})}^{+}},\quad
\Gamma_{\mathbb Q}=\mathrm{Gal}(\overline{\mathbb Q}_0/\mathbb Q)
\tag{9}
\]
the group maps (we presume injectively) into $\widehat{\widehat{\mathcal G}}(\pi_{g,\nu})$, \deleted{\unsure{by}} inducing an isomorphism between its subgroup $\pi_1(\mathcal U_{g,\nu\,\overline{\mathbb Q}_0})$ with $\widehat{\mathcal G}(\pi_{g,\nu})$, let us denote its image by $M_{g,\nu}$. So it is a \unsure{subgroup} \inserted{closed}
\[
\widehat{\mathcal G}_{g,\nu}\subset M_{g,\nu}\subset\widehat{\widehat{\mathcal G}}_{g,\nu}
\tag{10}
\]
and $\widehat{\mathcal G}_{g,\nu}^{+}$ is normal in $M_{g,\nu}$\inserted{ (but not $\widehat{\mathcal G}_{g,\nu}$!)}. The datum \deleted{\unsure{of}}\inserted{of such an} $M_{g,\nu}$ is equivalent to that of an $N_{g,\nu}$
\[
\widehat{\mathcal T}_{g,\nu}\subset N_{g,\nu}\subset\widehat{\widehat{\mathcal T}}_{g,\nu},
\tag{11}
\]
with $\widehat{\mathcal T}_{g,\nu}^{+}$ normal in $N_{g,\nu}$.

We set
\[
\Gamma_{g,\nu}=N_{g,\nu}/\widehat{\mathcal T}_{g,\nu}^{+}\simeq M_{g,\nu}/\widehat{\mathcal G}_{g,\nu}^{+}.
\tag{12}
\]
\margin{\uline{N.B.} In $\Gamma_{g,\nu}$ we have a canonical element of order $2$ $\tau_{g,\nu}$, \unsure{corresponding} to the elements of $\deleted{\unsure{A}}\widehat{\mathcal T}_{g,\nu}\unsure{-|\smallsetminus}\widehat{\mathcal T}_{g,\nu}^{+}=\widehat{\mathcal T}_{g,\nu}^{-}$.}

We have a surjective homomorphism
\[
\Gamma_{\mathbb Q}\longrightarrow\Gamma_{g,\nu}
\tag{13}
\]
which we presume is bijective --- i.e. that the $\Gamma_{g,\nu}$ are canonically isomorphic to one another.

\margin{370}

\scanpage{372}

Let $\pi_0$ now be a discrete group with loops of type $g,\nu$, then we define subgroups $M_{\inserted{\pi_0}}\deleted{(\unsure{\pi_0})}$, $N_{\inserted{\pi_0}}\deleted{(\unsure{\pi_0})}$
\[
\widehat{\mathcal G}(\pi_0)\subset M_{\inserted{\pi_0}}\deleted{(\unsure{\pi_0})}\subset\widehat{\widehat{\mathcal G}}(\pi_0)
\tag{14}
\]
\margin{$\Gamma_{\pi_0}=M_{\pi_0}/\widehat{\mathcal G}(\pi_0)$}
\[
\tikz[baseline=(lmfixA.base)]\node[draw,rounded corners=0.7em,inner sep=1.2pt](lmfixA){$\widehat{\widehat{\mathcal T}}$};(\pi_0)\subset N_{\inserted{\pi_0}}\deleted{(\unsure{\pi_0})}=M(\pi_0)/\widehat{\pi}_0\subset\widehat{\widehat{\mathcal T}}(\pi_0)
\tag{15}
\]
\margin{$\longleftrightarrow$\,? $\widehat{\mathcal T}$}
using an isomorphism $\pi_0\overset{\sim}{\longrightarrow}\pi_{g,\nu}$ and \unsure{proceeding} by transport of structure --- the result does not depend on it. More generally, let us start with a pair $(\pi,\alpha)$ of a profinite group $\pi$, equipped with a \uline{prediscretification} $\alpha$, which we can therefore interpret as a class of isomorphisms $\widehat{\pi}_{g,\nu}\longrightarrow\pi$, \unsure{defined|defined} modulo composition on the right by a $u\in\widehat{\mathcal G}_{g,\nu}$. Then by transport of structure we deduce from it groups \inserted{$M_\alpha$, $N_\alpha$}
\[
\widehat{\mathcal G}_{\inserted{\alpha}}\deleted{(\unsure{\pi})}\subset M_\alpha\subset\widehat{\widehat{\mathcal G}}(\pi)
\tag{16}
\]
\[
\widehat{\mathcal T}_{\inserted{\alpha}}\deleted{(\unsure{\pi})}\subset
\underset{\inserted{\substack{\Vert\\ M_\alpha/\pi}}}{N_\alpha}
\subset\tikz[baseline=(lmfixB.base)]\node[draw,rounded corners=0.7em,inner sep=1.2pt](lmfixB){$\widehat{\mathcal T}(\pi)$};
\tag{17}
\]
\margin{$\unsure{\widehat{\widehat{\mathcal T}}}$ ?}
with $\widehat{\mathcal G}(\pi)^{+}$ normal in $M_\alpha$ i.e. $\widehat{\mathcal T}^{+}(\pi)$ normal in $N_\alpha$. We set
\[
\Gamma_\alpha=M_\alpha/\widehat{\mathcal G}_{\inserted{\alpha}}^{\inserted{+}}\deleted{(\unsure{\pi})^{+}}\simeq N_\alpha/\widehat{\mathcal T}_\alpha^{+}
\tag{18}
\]
We have a canonical element \deleted{$\tau_\alpha$ \unsure{whose}}
\[
\tau_\alpha\in\Gamma_\alpha,\qquad \tau_\alpha^2=1,\quad\tau_\alpha\ne1
\tag{19}
\]

\margin{371}

\scanpage{373}
\mspage{189}

Let $\pi$ now be a profinite group with loops of type $(g,\nu)$. We call an \uline{\unsure{arithmetization}} of $\pi$ the datum of an element of
\[
\begin{aligned}
 A(\pi)&=\uline{\mathrm{Isom}}_{\mathrm{lac}}(\widehat{\pi}_{g,\nu},\pi)/M_{g,\nu}\\
 &\simeq\uline{\unsure{\mathrm{Isomext}}}_{\mathrm{lac}}(\widehat{\pi}_{g,\nu},\pi)/N_{g,\nu}
\end{aligned}
\tag{20}
\]
Let $P(\pi)$ be the \unsure{set} of prediscretifications \inserted{\unsure{oriented}} of $\pi$
\[
\begin{aligned}
 P(\pi)&=\uline{\mathrm{Isom}}_{\mathrm{lac}}(\widehat{\pi}_{g,\nu},\pi)/\widehat{\mathcal G}_{g,\nu}^{+}\\
 &\simeq\unsure{\mathrm{Isomext}}_{\mathrm{lac}}(\widehat{\pi}_{g,\nu},\pi)/\widehat{\mathcal T}_{g,\nu}^{+}
\end{aligned}
\tag{21}
\]
then $\Gamma_{g,\nu}$ operates freely \inserted{on the right} on $P(\pi)$, and
\[
 A(\pi)\simeq P(\pi)/\Gamma_{g,\nu}
\tag{22}
\]
--- $P(\pi)$ is a relative torsor over $A(\pi)$, with group $\Gamma_{g,\nu}$. For $a\in A(\pi)$, let $M_a$ be the subgroup of $\widehat{\widehat{\mathcal G}}(\pi)$ of \unsure{automorphisms} of $(\pi,a)$, it operates on the $\Gamma_{g,\nu}$-torsor $P_a$ of oriented discretifications of $\pi$ over $a$.
\inserted{For \deleted{\unsure{every}} $\alpha\in P_a$, let $\widehat{\mathcal G}_{\alpha}^{\unsure{*|\deleted{+}}}\subset\widehat{\widehat{\mathcal G}}(\pi)$, then $\widehat{\mathcal G}_{\alpha}^{+}$}
\deleted{is the subgroup of $M_a$ operating \unsure{trivi}}
\unsure{does} not depend on $\alpha$ (\unsure{we shall denote it|in view of the notation} (\illegible{1}\unsure{)!|?!} $\widehat{\mathcal G}_a$), it is also the kernel of the operation of $M_a$ on $P_a$. We have
\margin{\unsure{but. Careful} $\widehat{\mathcal G}_{\alpha}^{\unsure{*|\deleted{+}}}$ depends on $\alpha$ i.e. $\tau_{\alpha}\in\Gamma_a$ depends on $\alpha$}
\[
 \widehat{\mathcal G}_a^{+}\subset M_a\subset\widehat{\widehat{\mathcal G}}(\pi)
\tag{23}
\]
hence also
\margin{372}

\scanpage{374}

\[
 \widehat{\mathcal T}_a^{+}\subset N_a\subset\widehat{\widehat{\mathcal T}}(\pi)
\tag{24}
\]
setting
\[
 N_a\unsure{=}M_a/\pi,\qquad
 \widehat{\mathcal T}_a^{+}=\widehat{\mathcal G}_a^{+}/\pi.
\tag{25}
\]
We set
\[
 \Gamma_a=M_a/\widehat{\mathcal G}_a^{+}\simeq N_a/\widehat{\mathcal T}_a^{+}
\tag{26}
\]
then $\Gamma_a$ identifies with the commutant of $\Gamma_{g,\nu}$ operating on $P_a$, i.e. with $\Gamma_{g,\nu}$ ``twisted'' by the torsor $P_a$. The choice of an $\alpha\in P_a$ \unsure{amounts} therefore to the choice of such an \unsure{isomorphism}
\deleted{$\Gamma_a\unsure{\simeq}$}
\[
 \Gamma_a\overset{\sim}{\longleftrightarrow}\Gamma_{g,\nu}
\tag{27}
\]
which is \unsure{changed} by an inner \unsure{automorphism} if we change \unsure{$\alpha$|$\alpha$}...

We have the category of \deleted{\unsure{arith|ext}} arithmetized outer groups with loops of type $g,\nu$ --- it is a connected groupoid, with a fixed origin $(\widehat{\pi}_{g,\nu},\underbrace{a_{g,\nu}}_{\inserted{\text{\unsure{arithmetization} ``canonical''}}})$,
\deleted{--- in fact; the $\widehat{\pi}_0$ \unsure{corresponding} to discrete groups are canonically isomorphic \unsure{in|in} this category,}
\inserted{whose group of \unsure{automorphisms} is $N_{g,\nu}$, an \unsure{extension} of $\Gamma_{g,\nu}$ by $\widehat{\mathcal T}_{g,\nu}$, $\simeq\pi_1(M_{g,\nu\,\unsure{\mathbb Q|\overline{\mathbb Q}}})$ (\unsure{computed} with \unsure{respect} to the canonical universal covering of $M_{g,\nu\,\mathbb Q}$...)}
\margin{if $\Gamma_{\mathbb Q}\longrightarrow\Gamma_{g,\nu}$ injective!}
\margin{373}

\scanpage{375}
\mspage{190}

To give oneself an object of this category is essentially the $=$ thing (up to unique isomorphism) as to give oneself a universal covering \inserted{$\widetilde M_{g,\nu\,\mathbb Q}$} of $M_{g,\nu\,\unsure{\overline{\mathbb Q}}}$\inserted{$\unsure{\mathbb Q}$} --- or a torsor under $N_{g,\nu}$; we then \unsure{consider} the family of curves of type $g,\nu$ over $\unsure{\widetilde M_{g,\nu\,\mathbb Q}|M_{g,\nu\,\mathbb Q}}$
\[
 \mathcal U_{g,\nu}\times_{M_{g,\nu}}\widetilde M_{g,\nu\,\mathbb Q}
 =\mathcal U_{g,\nu}(\widetilde M_{g,\nu\,\mathbb Q})
\tag{28}
\]
as being ``the'' algebraic curve whose $\pi_1$ \inserted{outer} (which is \unsure{therefore} the $\pi_1$ of this scheme...) is the given outer group with loops.

If we took the arithmetized groups with loops, not outer, we would have \inserted{again} a connected groupoid with ``origin'' $(\widehat{\pi}_{g,\nu},a_{g,\nu})$, with a group of \unsure{automorphisms} which is $M_{g,\nu}$\inserted{$\simeq\pi_1(\mathcal U_{g,\nu\,\mathbb Q})$}. The datum of an object of this category \deleted{is} \inserted{amounts to} the datum of a universal covering (not only of $M_{g,\nu\,\mathbb Q}$, but) of $\mathcal U_{g,\nu\,\mathbb Q}$, i.e. $\widetilde{\mathcal U}_{g,\nu\,\mathbb Q}$, which indeed makes it possible to serve as a reference universal covering
\margin{374}

\scanpage{376}

for the \unsure{relative} curve (28), \unsure{allowing} one then to specify its fundamental group as a genuine group with \unsure{loops} (not just an outer group).

We can finally also look at the category of profinite groups with \unsure{loops} \unsure{arithmetized}, where we take as morphisms the \uline{arithmetically outer isomorphisms}.

We again find a connected groupoid with \unsure{marked} origin $(\pi_{g,\nu},a_{g,\nu})$, whose group of automorphisms is now $\Gamma_{g,\nu}$. Now the groups $\widehat{\pi}_0$ ($\pi_0$ a discrete group with \unsure{loops} of type $g,\nu$) are canonically isomorphic to one another. The category is (\unsure{conjecturally}, \unsure{assuming} $\Gamma_{\mathbb Q}\to\Gamma_{g,\nu}$ an isomorphism) equivalent to that of universal coverings of $\mathrm{Spec}\ \mathbb Q$, i.e. to that of algebraic closures of $\mathbb Q$, the origin element \unsure{corresponding} to $\overline{\mathbb Q}_0$.

\margin{375}

\scanpage{377}
\mspage{191}

When we give ourselves a $\pi$ with arithmetization $a$, then \unsure{there} \deleted{\unsure{gives|given}} corresponds to it \unsure{therefore} canonically an algebraic closure $\overline{\mathbb Q}$ of $\mathbb Q$. The \inserted{$\Gamma_{\mathbb Q}$} torsor of isomorphisms $\mathrm{Isom}(\overline{\mathbb Q}_0,\overline{\mathbb Q})$, i.e. of embeddings $\overline{\mathbb Q}\unsure{\hookrightarrow|\subset}\mathbb C$, is \unsure{in} $1-1$ correspondence with the set \inserted{$P_a$} of \unsure{oriented} prediscretifications \inserted{of $\pi$} \unsure{giving} rise to $a$.

Let us note that every $\alpha$ defines an element $\tau_\alpha\in\Gamma_a$, $\tau_\alpha^2=1$ ($\tau_\alpha\ne1$), hence a canonical map
\[
\tag{29}
P_a\longrightarrow {}_2\Gamma_a
\]
We see that this map is compatible with the action of the group $\pm1$ on $P_a$
\[
\tau_\alpha=\tau_{-\alpha}
\]

If it is true that $\Gamma_{\mathbb Q}\stackrel{\sim}{\longrightarrow}\Gamma_{g,\nu}$, then \deleted{the image \unsure{of} (29) is} (29) induces by passage to the quotient a \unsure{bijective} map
\[
\tag{30}
\begin{array}{ccl}
P_a/\pm1&\stackrel{\sim}{\longrightarrow}&\begin{gathered}\text{set of elements}\\\text{of order 2 of }\Gamma_a\end{gathered}\\
\inserted{\Vert}&&\\
\inserted{P_a^{\natural}=\begin{gathered}\text{set of prediscretifications}\\\text{(not oriented) over }a\unsure{)}\end{gathered}}&&
\end{array}
\]

\margin{376}

\scanpage{378}

This \unsure{comes} from the known fact that in $\Gamma_{\mathbb Q}$, the only elements of order 2 are the conjugates of $\tau$, and that the \unsure{centralizer} of $\tau$ in $\Gamma_{\mathbb Q}$ is reduced to $\{1,\tau\}$.

We can say that the datum of a \unsure{discretification} (not oriented) $\alpha^{\natural}$ over the arithmetization $a$, amounts to the datum of an \deleted{topology} Archimedean valuation on the algebraic closure $\overline{\mathbb Q}$ \inserted{of $\mathbb Q$} defined by $(\pi,a)$ (i.e. of an isomorphism $\overline{\mathbb Q}_1\simeq\overline{\mathbb Q}_0$ \uline{modulo complex conjugation})

\deleted{\unsure{30}}

\margin{377}

\scanpage{379}
\mspage{192}

\uline{30}\quad \uline{Properties of the} $N_{g,\nu}$, $\Gamma_{g,\nu}$.
\margin{\unsure{a)} Properties related to finite \unsure{subgroups|subgroups}.}

We would like to bring out properties, inspired by \unsure{the} geometrico-arithmetic context, but which \unsure{can|\deleted{can}} be \unsure{formulated|formulated} \unsure{in} a purely algebraic way -- and which \unsure{are} strong enough perhaps to end up \unsure{characterizing} the $N_{g,\nu}$, \deleted{\illegible{1}}

Let us return to a $(\pi,a)$, an arithmetized group with \unsure{loops}, heuristically, it corresponds to the datum of an algebraic closure $\overline{\mathbb Q}$ of $\mathbb Q$, and (if $\pi$ is given as an outer group) of an algebraic family, \unsure{parametrized} by a universal covering \inserted{$\unsure{\widetilde M_{g,\nu\overline{\mathbb Q}}}$} of $M_{g,\nu\overline{\mathbb Q}}$, of algebraic curves of type $(g,\nu)$. In this perspective, to reduce this family to \uline{one} algebraic curve -- i.e. to give oneself an algebraic curve of type $(g,\nu)$ over $\overline{\mathbb Q}$ -- must amount to giving oneself a kernel of lifting of the surjective homomorphism
\[
 N_a\longrightarrow\Gamma_a
 \qquad (\text{kernel }\widehat{\mathcal T}_a^+)
 \tag{1}
\]
\margin{378}

\scanpage{380}

The datum of such a lifting over an open subgroup $\Gamma'$, amounts to the datum of an algebraic curve of type $g,\nu$, defined over \tikz[baseline=(lmfixC.base)]\node[draw,rounded corners=0.7em,inner sep=1.2pt](lmfixC){\shortstack{\scriptsize\unsure{hum-|hom-}\\the extension}}; finite $K'\subset\overline{\mathbb Q}$ defined by $\Gamma'$.
\margin{?}

It is a matter here not \unsure{of|of the} arbitrary continuous liftings, but of liftings having particular properties (some of which very deep, of the ``Weil'' kind\ldots\deleted{)} but \uline{perhaps} consequences of much simpler properties). The properties \unsure{\deleted{thus} to be brought out|\deleted{thus brought out}} \unsure{should} in any case be \deleted{\unsure{invariant}} \inserted{stable} under passage to a smaller open subgroup. Among these properties, there would be that \unsure{these|a} \unsure{germs|germ} of liftings \unsure{could|could} in their turn be lifted to $\widehat{\widehat{\mathcal G}}(\pi)$ itself -- \unsure{in} an equally ``admissible'' way in a sense to be specified -- and \unsure{=|even} that there would be ``many'' \inserted{(\unsure{classes} of $\pi$-conjugacy)} of such
\margin{379}